%! Transformations of Functions — Unit 2, Lesson 2.2 — AP Practice
% EXTRA CREDIT — optional, exactly TWO pages, placed between the notes and the graded homework.
% Page 1 is Section I, five multiple-choice items with four options (A)–(D); page 2 is Section
% II, one multi-part free-response question on a table, in context. The greenhouse context is
% used nowhere else in the lesson.
% Distractors are the lesson's errors on purpose: MC 1 reads the inside sign forwards; MC 2
% reads f(3x) as a stretch by 3 (and applies the 3 to y); MC 3 moves the asymptote the wrong
% way; MC 4 flips the inside sign or the outside one, or inverts a; MC 5 forgets that a
% reflection flips the range (EK 1.12.A.6). Answers: C, D, B, A, C.
%
% PAGE LOCKSTEP with the other half of the pair: the key marks the correct option in its
% LABEL (no text added to the line), the table blanks become \ans, and prose answers are
% \writespace{H}{…}, fixed height both ways. The explicit \newpage fixes the page break.
\documentclass[10pt]{article}
\usepackage{precalculus-article}
\usepackage{precalculus-boxes}
\newcommand{\TallMath}[1]{$\displaystyle #1\rule[-1.4em]{0pt}{3.2em}$}

\begin{document}

\pageheader{Unit 2, Lesson 2.2}{AP Practice}

\vspace{0.12in}

\boxguard[8]
\begin{remindbox}
\small
\textbf{Extra credit --- optional.} These two pages are written in the style of the AP
Precalculus exam. Your graded homework is the last section of this packet; do it first. Every
answer choice below is a mistake someone really makes, so read all four before you choose.
\end{remindbox}

\vspace{0.12in}

\begin{headlinebox}{goldbg}
\small
\textbf{Section I --- Multiple Choice.} Circle the letter of the single best answer. No calculator.
\end{headlinebox}

\vspace{0.06in}

\begin{enumerate}[leftmargin=1.9em, itemsep=12pt]

\item Which function's graph is the graph of $y = \sqrt{x}$ translated 2 units to the
      \textbf{left} and 5 units \textbf{down}?

      \begin{enumerate}[leftmargin=2.6em, itemsep=2pt]
        \item[(A)] $y = \sqrt{x - 2} - 5$
        \item[(B)] $y = \sqrt{x - 5} + 2$
        \item[(C)] $y = \sqrt{x + 2} - 5$
        \item[(D)] $y = \sqrt{x + 5} - 2$
      \end{enumerate}

\item The point $(12, -6)$ lies on the graph of the function $f$. Let $g(x) = f(3x)$. Which
      point must lie on the graph of $g$?

      \smallskip
      \textbf{(A)} $(36, -6)$ \hspace{2.2em} \textbf{(B)} $(12, -18)$ \hspace{2.2em} \textbf{(C)} $(4, -2)$ \hspace{2.2em} \textbf{(D)} $(4, -6)$

\item Let $g(x) = \dfrac{1}{x - 2} + 3$, a transformation of $y = \dfrac{1}{x}$. Which gives the
      vertical and horizontal asymptotes of the graph of $g$?
      \begin{enumerate}[leftmargin=2.6em, itemsep=2pt]
        \item[(A)] $x = -2$ and $y = 3$
        \item[(B)] $x = 2$ and $y = 3$
        \item[(C)] $x = 3$ and $y = 2$
        \item[(D)] $x = 2$ and $y = -3$
      \end{enumerate}

\item \begin{minipage}[t]{0.56\linewidth}
      The graph of a transformation of $y = x^2$ is shown. Its vertex is $(1, -2)$, and it
      passes through $(0, -4)$ and $(2, -4)$. Which is an equation for the graph?
      \begin{enumerate}[leftmargin=2.6em, itemsep=3pt]
        \item[(A)] $y = -2(x - 1)^2 - 2$
        \item[(B)] $y = -2(x + 1)^2 - 2$
        \item[(C)] $y = 2(x - 1)^2 - 2$
        \item[(D)] $y = -\tfrac{1}{2}(x - 1)^2 - 2$
      \end{enumerate}
      \end{minipage}\hfill
      \begin{minipage}[t]{0.38\linewidth}
      \vspace{-4pt}
      \centering
      \begin{tikzpicture}[x=0.5cm, y=0.3cm, baseline=(current bounding box.north)]
        \draw[linegray, very thin, step=1] (-2,-9) grid (4,2);
        \draw[->, thick] (-2,0) -- (4.5,0) node[right] {\scriptsize $x$};
        \draw[->, thick] (0,-9) -- (0,2.5) node[above] {\scriptsize $y$};
        \foreach \x in {-1,1,2,3} \node[below, font=\tiny] at (\x,0) {\x};
        \foreach \y in {-8,-6,-4,-2} \node[left, font=\tiny] at (0,\y) {\y};
        \draw[very thick, plum, domain=-1:3, samples=50] plot (\x, {-2*(\x-1)*(\x-1) - 2});
        \fill[plum] (1,-2) circle (2.2pt);
        \fill[plum] (0,-4) circle (2.0pt);
        \fill[plum] (2,-4) circle (2.0pt);
      \end{tikzpicture}
      \end{minipage}

\item The range of $f(x) = \sqrt{x}$ is $y \ge 0$. What is the range of
      $g(x) = -3\sqrt{x} + 4$?

      \smallskip
      \textbf{(A)} $y \ge 4$ \hspace{2.2em} \textbf{(B)} $y \ge 0$ \hspace{2.2em} \textbf{(C)} $y \le 4$ \hspace{2.2em} \textbf{(D)} $y \le -3$

\end{enumerate}

\newpage

\begin{headlinebox}{goldbg}
\small
\textbf{Section II --- Free Response.} Show your reasoning. Answers about the greenhouses
must be written \emph{in context}, with units --- that is what earns credit on the AP exam.
\end{headlinebox}

\vspace{0.08in}

\begin{enumerate}[leftmargin=1.9em, itemsep=10pt, start=6]

\item The temperature in Greenhouse A, in degrees Celsius, is $f(t)$, where $t$ is hours after
      noon (a negative $t$ is before noon). A few values are given. Greenhouse B runs on a
      different timer and heater, and its temperature is $g(t) = 2f(t - 3) + 1$.

      \begin{center}
      {\renewcommand{\arraystretch}{1.3}%
      \begin{tabular}{|c|c|c|c|c|}
      \hline
      $t$ (hours) & $-2$ & $0$ & $2$ & $4$ \\ \hline
      $f(t)$ ($^\circ$C) & $3$ & $-1$ & $1$ & $5$ \\ \hline
      \end{tabular}}
      \end{center}

      \textbf{(a)} Each point $(t, f(t))$ in the table becomes a point of $g$. Complete the
      table.

      \begin{center}
      {\renewcommand{\arraystretch}{1.35}%
      \begin{tabular}{|c|c|c|c|c|}
      \hline
      Point of $f$ & $(-2, 3)$ & $(0, -1)$ & $(2, 1)$ & $(4, 5)$ \\ \hline
      Point of $g$ & \blank{1.4cm} & \blank{1.4cm} & \blank{1.4cm} & \blank{1.4cm} \\ \hline
      \end{tabular}}
      \end{center}

      \textbf{(b)} Find $g(5)$, and name the value of $f$ it used. Interpret $g(5)$ in the
      context of the problem, with units.
      \writespace{2.0cm}{}

      \textbf{(c)} Describe, in order, the transformations that turn the graph of $f$ into the
      graph of $g$.
      \writespace{2.0cm}{}

      \textbf{(d)} A third greenhouse runs on a faster cycle, $h(t) = f(2t)$. Kai claims the
      graph of $h$ passes through $(8, 5)$, because $f(4) = 5$. Is Kai right? Find the correct
      point and explain.
      \writespace{2.4cm}{}

      \textbf{(e)} Write a function $j$, in terms of $f$, whose graph is the graph of $f$ reflected
      over the $y$-axis and then moved up 4. Which point of $j$ comes from $(4, 5)$?
      \writespace{2.2cm}{}

\end{enumerate}

\end{document}
